\documentclass[aps,pre,twocolumn,superscriptaddress,floatfix]{revtex4-2}

\usepackage{amsmath,amssymb}
\usepackage{graphicx}
\usepackage{booktabs}
\usepackage{hyperref}

\begin{document}

\title{Numerical Hessian Rank and Keplerian Eigenvalue Structure\\as Diagnostics for Periodic Orbits in the Equal-Mass Three-Body Problem}

\author{Andrew H. Bond}
\email{andrew.bond@sjsu.edu}
\affiliation{Department of Computer Engineering, San Jose State University, San Jose, CA 95192, USA}
\affiliation{Senior Member, IEEE; ORCID: 0009-0003-2599-6158}

\date{\today}

\begin{abstract}
We compute the Hessian of the gravitational three-body Hamiltonian in
Jacobi coordinates across 3.24 million configurations and analyze its
numerical rank, eigenvalue spectrum, and relationship to orbital
periodicity. At configurations where the numerical rank of the
potential-energy block $V_{qq}$ drops to~3 (using a singular-value
threshold $\epsilon = 10^{-6}$), the three nonzero eigenvalues
exhibit an exact $2{:}1{:}1$ magnitude ratio---the curvature signature
of a dominant Keplerian binary with a decoupled third body. This ratio
is proven analytically from the Hessian of the $1/r$ potential and
confirmed computationally to $< 10^{-12}$ relative error. We integrate
100{,}000 orbits on GPU and find that configurations at effective
rank~9 (potential rank~3) show a periodicity rate of $5.5\%$,
the highest among all rank levels.
A gradient-boosted classifier trained on Hessian-derived features
achieves cross-validated $F_1 = 0.86$ for periodicity prediction on
a 51{,}000-configuration dataset, and an A*-discovered closed-form
formula $\sqrt{n_\mathrm{clusters}^2 + \lambda_\mathrm{min}^2} > \tau$
achieves $F_1 = 0.48$. The rank classification is threshold-dependent
(vanishing at $\epsilon = 10^{-10}$) and the Keplerian structure
is asymptotic (exact only in the decoupled limit), which we
characterize through threshold sensitivity sweeps. We compare against
separation ratio as a baseline and find that Hessian features provide
complementary but not superior predictive power for periodicity.
\end{abstract}

\maketitle

\section{Introduction}

The gravitational three-body problem has no general closed-form
solution~\cite{poincare1892}, but specific configurations---hierarchical
triples, Lagrange points, choreographies---are approximately or exactly
tractable. Identifying which configurations are tractable typically
requires either perturbation theory~\cite{murray1999solar}, KAM
analysis~\cite{arnold1963small}, or brute-force numerical
integration~\cite{suvakov2013three}. Each approach is computationally
expensive or requires analytical insight about the specific
configuration.

We explore whether the \emph{Hessian} of the Hamiltonian---computed
once per configuration at negligible cost compared to orbit
integration---contains sufficient information to predict orbital
regularity. The Hessian encodes the local curvature of the energy
landscape: its rank measures how many degrees of freedom are
dynamically coupled, and its eigenvalue structure reflects the
characteristic frequencies of the linearized motion.

Our approach is computational rather than theoretical. We make no
claim that Hessian rank is a superior diagnostic to existing methods;
rather, we characterize what it reveals and where it fails, as a
contribution to the toolkit of three-body diagnostics.

\section{Methods}

\subsection{Jacobi Coordinates and Hessian Structure}

We work in the center-of-mass-reduced 12-dimensional phase space
$z = (\boldsymbol{\rho}_1, \boldsymbol{\rho}_2, \boldsymbol{\pi}_1,
\boldsymbol{\pi}_2)$. The Hamiltonian separates as
$H = T(\boldsymbol{\pi}) + V(\boldsymbol{\rho})$ with
$T = |\boldsymbol{\pi}_1|^2/(2\mu_1) + |\boldsymbol{\pi}_2|^2/(2\mu_2)$.

The Hessian $C_{ab} = \partial^2 H / \partial z_a \partial z_b$ has
exact block-diagonal structure:
\begin{equation}
C = \begin{pmatrix} V_{qq} & 0 \\ 0 & T_{pp} \end{pmatrix},
\end{equation}
where the position-momentum cross block is identically zero (since
$H$ is separable in $q$ and $p$), and $T_{pp} = \mathrm{diag}(\mu_1^{-1}
I_3, \mu_2^{-1} I_3)$ is configuration-independent with eigenvalues
$\{1/\mu_1, 1/\mu_1, 1/\mu_1, 1/\mu_2, 1/\mu_2, 1/\mu_2\}$.
For equal masses, these are $\{2, 2, 2, 1.5, 1.5, 1.5\}$.

All variation in the Hessian structure comes from the $6 \times 6$
potential block $V_{qq}$, which we compute analytically (see
Appendix). The total Hessian rank satisfies
$\mathrm{rank}(C) = \mathrm{rank}(V_{qq}) + 6$.

\subsection{Numerical Rank}

We define the \emph{effective rank} at threshold $\epsilon$ as:
\begin{equation}
r_\epsilon(A) = |\{k : \sigma_k / \sigma_1 > \epsilon\}|
\end{equation}
where $\sigma_k$ are the singular values of $A$ in descending order.
Unless stated otherwise, we use $\epsilon = 10^{-6}$.

This is a numerical, not algebraic, quantity. At finite body
separations, $V_{qq}$ is generically full rank (all six eigenvalues
nonzero). The ``rank~3'' designation means that three singular values
are below $10^{-6}$ times the largest---not that they are exactly zero.
Table~\ref{tab:threshold} shows the sensitivity to $\epsilon$.

\begin{table}[h]
\caption{Fraction of configurations with effective rank $< 12$
at various thresholds ($N = 10{,}000$ random sample, equal masses).}
\label{tab:threshold}
\begin{ruledtabular}
\begin{tabular}{cc}
Threshold $\epsilon$ & Fraction rank $< 12$ \\
\hline
$10^{-4}$ & 44.8\% \\
$10^{-5}$ & 25.8\% \\
$10^{-6}$ (default) & 8.2\% \\
$10^{-7}$ & 2.9\% \\
$10^{-8}$ & 1.0\% \\
$10^{-10}$ & 0.0\% \\
\end{tabular}
\end{ruledtabular}
\end{table}

\subsection{Configuration Sampling}

We sampled $50 \times 50 \times 36 \times 36 = 3{,}240{,}000$
configurations for equal masses ($m_1 = m_2 = m_3 = 1$) on a grid
in $(r_1, r_2, \theta, \phi)$ with logarithmic radial spacing
$r \in [0.1, 100]$, where $r_1$ is the inner pair separation and
$r_2$ is the distance from the third body to the inner pair's
center of mass. Initial momenta are set for circular orbits.

\subsection{Orbit Integration}

Orbits are integrated using a symplectic leapfrog (Störmer-Verlet)
scheme with $\Delta t = 0.005$ for $8 \times 10^4$ steps ($T = 400$
natural time units). An orbit is classified as \emph{periodic} if it
returns to within 5\% of its initial phase-space extent at any time
after the first 10\% of the integration. Energy conservation error
is monitored; mean relative error is $\sim 10^{-3}$.

Integration is GPU-accelerated using CuPy on an NVIDIA Quadro GV100
(32~GB), processing up to 400{,}000 orbits simultaneously. The
maximum batch size is determined automatically using the
batch-probe library~\cite{batchprobe}.

\section{Results}

\subsection{Rank Landscape}

Table~\ref{tab:landscape} shows the complete rank distribution.

\begin{table}[h]
\caption{Hessian rank distribution (equal masses, $\epsilon = 10^{-6}$,
$N = 3{,}240{,}000$).}
\label{tab:landscape}
\begin{ruledtabular}
\begin{tabular}{crrc}
Rank & Count & Fraction & $V_{qq}$ eff.\ rank \\
\hline
6 & 24 & 0.001\% & 0 \\
7 & 1{,}316 & 0.04\% & 1 \\
8 & 716 & 0.02\% & 2 \\
9 & 156{,}756 & 4.84\% & 3 \\
10 & 74{,}124 & 2.29\% & 4 \\
11 & 32{,}972 & 1.02\% & 5 \\
12 & 2{,}974{,}092 & 91.79\% & 6 \\
\end{tabular}
\end{ruledtabular}
\end{table}

\subsection{Keplerian Eigenvalue Structure}

At configurations where $V_{qq}$ has effective rank~3, the three
dominant eigenvalues exhibit a $2{:}1{:}1$ magnitude ratio:
\begin{equation}
|\lambda_1| / |\lambda_2| = |\lambda_1| / |\lambda_3| = 2.000\,000
\label{eq:ratio}
\end{equation}
to better than $10^{-6}$ relative error across all 156{,}756 such
configurations. The eigenvalue signs are $\{-\lambda, +\lambda/2,
+\lambda/2\}$ where $\lambda \approx 2m_1 m_2 / r_1^3$.

\emph{Proof.} In the limit where the third body is decoupled,
$V \approx -m_1 m_2 / |\boldsymbol{\rho}_1|$. The Hessian of
$-m/r$ in 3D evaluated at $\boldsymbol{\rho}_1 = (r_1, 0, 0)$ is:
\[
\frac{\partial^2}{\partial x_i \partial x_j}\!\left(-\frac{m}{r}\right)
= \frac{m}{r^3}\!\left(\frac{3 x_i x_j}{r^2} - \delta_{ij}\right)
\]
with eigenvalues $+2m/r^3$ (radial, multiplicity~1) and $-m/r^3$
(tangential, multiplicity~2). The magnitude ratio is exactly $2{:}1$.
At finite third-body distance, corrections of order
$(r_1/r_2)^3$ perturb the eigenvalues but preserve the ratio to high
accuracy when $r_2 \gg r_1$. $\square$

This result is the Keplerian curvature structure expressed in the
language of the coupling tensor. It is not new physics---it follows
directly from the $1/r$ potential---but it provides a computationally
cheap diagnostic: configurations exhibiting the $2{:}1$ ratio are
near-Keplerian hierarchical triples.

\subsection{Periodicity by Rank}

Table~\ref{tab:orbits} shows periodicity rates from the 100{,}000-orbit
GPU survey.

\begin{table}[h]
\caption{Periodicity rate by effective Hessian rank (100{,}000 orbits,
equal masses, circular initial momenta).}
\label{tab:orbits}
\begin{ruledtabular}
\begin{tabular}{crrc}
Rank & Periodic & Total & Rate \\
\hline
7--8 & 0 & 66 & 0.0\% \\
9 & 270 & 4{,}905 & 5.5\% \\
10 & 39 & 2{,}345 & 1.7\% \\
11 & 0 & 1{,}048 & 0.0\% \\
12 & 2{,}975 & 91{,}635 & 3.2\% \\
\end{tabular}
\end{ruledtabular}
\end{table}

Rank~9 has the highest periodicity rate (5.5\%), but rank~12
(generic, full rank) also shows 3.2\% periodicity. The rank
diagnostic enriches the periodic population by a factor of
$5.5/3.2 \approx 1.7\times$ relative to the full-rank baseline---a
modest but nonzero signal.

\subsection{Comparison to Separation Ratio Baseline}

The simplest diagnostic for hierarchical stability is the separation
ratio $r_2/r_1$. On the 51{,}000-configuration labeled dataset
(three mass ratios), we compare Hessian-based predictors to this
baseline:

\begin{table}[h]
\caption{Periodicity prediction: Hessian features vs.\ baseline.}
\label{tab:comparison}
\begin{ruledtabular}
\begin{tabular}{lc}
Predictor & F1 \\
\hline
Random (base rate 5--10\%) & $\sim$0.07 \\
$r_2/r_1 > 10$ & 0.29 \\
Frequency clusters $\geq 2$ & 0.29 \\
A*-discovered formula (Eq.~\ref{eq:formula}) & 0.48 \\
Decision tree (depth 5, 9 features) & 0.58 \\
Gradient boosting (100 trees) & 0.86 \\
\end{tabular}
\end{ruledtabular}
\end{table}

The A*-discovered formula:
\begin{equation}
\sqrt{n_\mathrm{clusters}^2 + \lambda_\mathrm{min}^2} > \tau
\label{eq:formula}
\end{equation}
achieves F1 $= 0.48$, outperforming the separation ratio baseline
(F1 $= 0.29$) but well below the gradient-boosted ensemble
(F1 $= 0.86$). The frequency cluster count alone matches the
separation ratio---it is a Hessian-derived equivalent of the same
information, not a new signal.

\subsection{Machine Learning Predictor}

The gradient-boosted classifier (F1 $= 0.86$, 5-fold cross-validated)
uses 9 Hessian-derived features. Feature importances reveal that
\emph{energy} (38\%), \emph{log separation ratio} (18\%), and
\emph{log frequency ratio} (12\%) dominate. Tensor rank contributes
only 10\%. The decision tree identifies the primary split as
rank~$\leq 9.5$ (importance 46\%), recovering the rank threshold
automatically.

The gap between the closed-form formula (F1 $= 0.48$) and the
ensemble (F1 $= 0.86$) quantifies the irreducible complexity of the
periodicity prediction task: 38\% of the predictive structure requires
nonlinear feature interactions that no depth-4 formula can express.

\section{Discussion}

The Hessian rank landscape provides a computationally cheap ($97\times$
faster than integration) map of the three-body configuration space
that correlates with but does not strongly predict orbital periodicity.
The Keplerian $2{:}1$ eigenvalue ratio is a clean analytical result
but amounts to detecting near-Keplerian hierarchical structure---information
also available from the separation ratio.

The principal value of the Hessian approach is in the machine learning
setting: the 9 Hessian-derived features, when combined nonlinearly,
achieve F1 $= 0.86$ for periodicity prediction, substantially
outperforming any single diagnostic. This suggests that the Hessian
encodes complementary information about the energy landscape curvature
that simple geometric diagnostics (like separation ratio) miss.

\subsection{Limitations}

The rank classification is threshold-dependent (Table~\ref{tab:threshold})
and should not be interpreted as exact algebraic rank. All ``rank~9''
configurations have six nonzero singular values; three are merely small.
The $2{:}1$ eigenvalue ratio is exact only in the fully decoupled limit
and approximate (with $(r_1/r_2)^3$ corrections) at finite separations.
The periodicity classification depends on integration parameters (timestep,
duration, return threshold). Results are shown for equal masses only;
generalization to unequal masses shows qualitatively similar patterns
but quantitatively different rank distributions. The 4-body generalization
(3{,}000 configurations, equal masses) found 0\% periodicity, indicating
that the approach does not trivially extend to $N > 3$.

\section{Conclusion}

The Hessian of the three-body Hamiltonian provides a fast, analytically
grounded diagnostic for near-Keplerian hierarchical structure via its
numerical rank and $2{:}1$ eigenvalue ratio. As a standalone periodicity
predictor, its power is modest (F1 $= 0.48$ for the best formula,
$1.7\times$ enrichment at rank~9). Combined with energy and separation
features in a machine learning ensemble, it contributes to a strong
predictor (F1 $= 0.86$). The rank landscape atlas (3.24 million
configurations) and 100{,}000-orbit GPU dataset are publicly available
for further analysis.

\begin{acknowledgments}
Computations performed on a dual-GV100 workstation using CuPy for
GPU-accelerated batch integration. Optimal GPU batch sizes determined
using the batch-probe library~\cite{batchprobe}.
Code and data: \url{https://github.com/ahb-sjsu/astar-symbolic-discovery}.
\end{acknowledgments}

\begin{thebibliography}{10}

\bibitem{poincare1892}
H.~Poincar{\'e}, \emph{Les m{\'e}thodes nouvelles de la
m{\'e}canique c{\'e}leste} (Gauthier-Villars, Paris, 1892).

\bibitem{murray1999solar}
C.~D.~Murray and S.~F.~Dermott, \emph{Solar System Dynamics}
(Cambridge University Press, 1999).

\bibitem{arnold1963small}
V.~I.~Arnold, Russ.\ Math.\ Surv.\ \textbf{18}, 85 (1963).

\bibitem{suvakov2013three}
M.~\v{S}uvakov and V.~Dmitra\v{s}inovi\'{c}, Phys.\ Rev.\ Lett.\
\textbf{110}, 114301 (2013).

\bibitem{batchprobe}
A.~H.~Bond, batch-probe: Find the maximum batch size that fits in
GPU memory, v0.3.0 (2026). \url{https://pypi.org/project/batch-probe/}

\end{thebibliography}

\end{document}
